% !TEX root = cheatsheet.tex
\begin{paracol}{2}
\section*{QA3b. Theoretical Q\&A — Segmentation, BraTS \& Papers}

\textbf{Q1. End-to-end pipeline from patient to clinical image.}\\
The full chain has \textbf{seven stages}: (1) \textbf{Acquisition} — patient scanned by CT/MRI/US/PET; raw \textbf{k-space} or \textbf{sinogram} data captured. (2) \textbf{Reconstruction} — inverse Fourier (MRI) or filtered back-projection (CT) builds a volume. (3) \textbf{Preprocessing} — \textbf{N4 bias correction}, denoising, \textbf{intensity normalization} (Nyul), \textbf{skull-stripping}. (4) \textbf{Registration} — align modalities or atlas to patient space. (5) \textbf{Segmentation} — extract organs/lesions (U-Net, nnU-Net). (6) \textbf{Quantification} — volume, shape, radiomics features. (7) \textbf{Reporting / CDSS} — visualization, DICOM-SR, decision support to clinician. Each step propagates error, so QC at every stage is critical.

\textbf{Q2. Bias field correction (N4) — why for MRI?}\\
\textbf{Bias field} is a smooth, low-frequency \textbf{multiplicative} intensity inhomogeneity caused by \textbf{B1 RF coil non-uniformity} in MRI; the same tissue (e.g., white matter) appears brighter on one side of the brain and darker on the other. \textbf{N4ITK} (Tustison 2010) is an improved \textbf{N3} algorithm that models the log-bias as a \textbf{B-spline} surface and iteratively sharpens the histogram via \textbf{deconvolution}. Without N4, intensity-based segmentation (Otsu, k-means, even CNNs) fails because the same tissue class has different intensities. N4 is mandatory in BraTS preprocessing. CT does not need it (Hounsfield units are calibrated).

\textbf{Q3. Image registration — rigid vs affine vs B-spline vs Demons vs LDDMM.}\\
\textbf{Registration} aligns image $I$ to reference $J$ by finding transform $T$ minimizing dissimilarity. \textbf{Rigid} (6 DoF: 3 rot + 3 trans) — same patient, same modality, different time. \textbf{Affine} (12 DoF) — adds scaling + shear; useful for inter-subject coarse alignment. \textbf{B-spline FFD} — non-rigid, control-grid deformation; smooth, invertibility not guaranteed. \textbf{Demons} (Thirion) — optical-flow style, fast, no explicit regularizer beyond Gaussian smoothing. \textbf{LDDMM} (Large Deformation Diffeomorphic Metric Mapping) — geodesic in space of diffeomorphisms; guarantees \textbf{topology preservation} (no folds), used in computational anatomy.

\textbf{Q4. Similarity metrics in registration.}\\
Choice depends on modality match. \textbf{SSD} (sum-of-squared-differences) $\sum (I-J)^2$ — only for same modality, same intensity scale. \textbf{NCC} (normalized cross-correlation) — robust to linear intensity changes, used in mono-modal MRI. \textbf{Mutual Information} \(\mathrm{MI}(I,J)=\sum p(i,j)\log\frac{p(i,j)}{p(i)p(j)}\) — gold standard for \textbf{multi-modal} (MRI$\leftrightarrow$CT, T1$\leftrightarrow$T2) since it only assumes statistical dependence, not intensity equality. \textbf{NMI} (normalized MI) reduces overlap bias. For deep-learning registration (VoxelMorph) we use \textbf{NCC + smoothness loss} on the displacement field.

\textbf{Q5. Common denoising methods for medical images.}\\
\textbf{Gaussian filter} — simple, blurs edges. \textbf{Median filter} — removes salt-and-pepper, preserves edges, used on US speckle. \textbf{Bilateral filter} — edge-preserving via spatial+intensity kernel. \textbf{Anisotropic diffusion} (Perona-Malik) — PDE that diffuses along but not across edges. \textbf{Non-Local Means (NLM)} — averages similar patches across the image; gold standard before deep learning. \textbf{BM3D} — block-matching + 3D collaborative filtering, near-optimal for Gaussian noise. \textbf{DnCNN / Noise2Noise} — deep learning, learns residual noise. For MRI \textbf{Rician noise} model is needed.

\textbf{Q6. Intensity normalization (Nyul) — why across scanners?}\\
MRI intensities are \textbf{not calibrated} (unlike CT Hounsfield units): the same brain on a Siemens 3T and a Philips 1.5T can have wildly different intensity ranges. This destroys generalization. \textbf{Nyul \& Udupa (2000)} normalization picks a set of \textbf{histogram landmarks} (e.g., percentiles 1, 10, 20, ..., 99) on a training corpus, computes their mean positions, then \textbf{piecewise-linearly maps} every new scan's landmarks to those targets. Alternatives: \textbf{z-score} per-volume, \textbf{white-stripe} (normalize to white-matter peak), \textbf{histogram matching}. Without normalization, a model trained on hospital A fails on hospital B — a major \textbf{domain-shift} problem.

\textbf{Q7. What is segmentation? Why most important step?}\\
\textbf{Segmentation} assigns a label to every voxel: $S:\Omega\to\{0,1,\ldots,K\}$, partitioning the image into anatomically/pathologically meaningful regions. It is the bridge from pixels to \textbf{quantitative biomarkers}: tumor \textbf{volume}, growth rate (RECIST), shape descriptors, radiomics. Diagnosis ("is there a tumor?"), treatment planning (radiotherapy contours), surgical guidance, and outcome prediction all require segmentation. Classification only says "yes/no"; segmentation says "where + how big". This is why BraTS, LiTS, KiTS challenges all center on segmentation.

\textbf{Q8. Otsu thresholding — formula.}\\
\textbf{Otsu (1979)} finds the threshold $t^*$ that \textbf{maximizes between-class variance} (equivalent to minimizing within-class variance) over the gray-level histogram. With class probabilities $\omega_0(t),\omega_1(t)$ and means $\mu_0(t),\mu_1(t)$:
\[t^*=\arg\max_t \;\sigma_B^2(t)=\omega_0(t)\,\omega_1(t)\,[\mu_0(t)-\mu_1(t)]^2.\]
It assumes a \textbf{bimodal histogram} (foreground + background). Pros: parameter-free, fast $O(L)$. Cons: fails on multi-modal histograms, sensitive to noise and uneven illumination (use \textbf{adaptive/local Otsu}). Multi-class extension exists (\textbf{multi-Otsu}).

\textbf{Q9. Region growing.}\\
\textbf{Region growing} starts from one or more \textbf{seed pixels} and iteratively absorbs neighbors satisfying a \textbf{homogeneity criterion} (e.g., $|I(p)-\mu_R|<\tau$ or gradient-based). It is \textbf{semi-automatic} (user picks seed) and intuitive, used in liver/lung segmentation in older PACS tools. Weaknesses: highly seed-sensitive, suffers from \textbf{leakage} when boundaries are weak, can't separate touching structures, and threshold $\tau$ is hard to tune. Variants: \textbf{symmetric region growing}, \textbf{statistical region merging}, \textbf{confidence connected} (ITK).

\textbf{Q10. Watershed — over-segmentation \& fix.}\\
\textbf{Watershed} treats the gradient image as a topographic surface; "flooding" from local minima creates basins separated by \textbf{watershed lines} along ridges. Each basin = one region. Problem: \textbf{every local minimum becomes a region} → severe \textbf{over-segmentation} on noisy gradients (hundreds of basins). Fixes: (1) \textbf{Marker-controlled watershed} — flood only from user/algorithm-supplied markers; (2) \textbf{H-minima transform} — suppress shallow minima; (3) \textbf{Pre-smoothing} the gradient (Gaussian/bilateral); (4) \textbf{Region merging} as post-process based on similarity. Marker-controlled watershed is the standard variant used in cell segmentation.

\textbf{Q11. Active contours / snakes — energy formula.}\\
A \textbf{snake} is a parametric curve $C(s)=(x(s),y(s)),\,s\in[0,1]$ that evolves to minimize:
\[E_{\text{snake}}=\int_0^1\!\!\big(\underbrace{\alpha|C'(s)|^2+\beta|C''(s)|^2}_{E_{\text{int}}}+\underbrace{E_{\text{ext}}(C(s))}_{-|\nabla I|^2}\big)\,ds.\]
\textbf{Internal energy}: $\alpha$ controls elasticity (resists stretching), $\beta$ controls rigidity (resists bending). \textbf{External energy} pulls the curve toward edges (high $|\nabla I|$). Solved by gradient descent with finite-difference Euler-Lagrange. Weaknesses: needs good initialization, can't change topology, fails in concavities. \textbf{GVF snake} (Gradient Vector Flow) extends capture range.

\textbf{Q12. Level set methods — evolution PDE.}\\
\textbf{Level set} (Osher-Sethian 1988) embeds the contour as the zero level set of a higher-dimensional function $\phi(\mathbf{x},t)$: $C(t)=\{\mathbf{x}:\phi(\mathbf{x},t)=0\}$. Evolution PDE:
\[\frac{\partial\phi}{\partial t}+F\,|\nabla\phi|=0,\]
where speed $F$ depends on curvature $\kappa=\nabla\!\cdot\!(\nabla\phi/|\nabla\phi|)$ and image features. \textbf{Chan-Vese} model uses region statistics (mean inside/outside) instead of edges. Big advantage over snakes: \textbf{automatic topology change} (split/merge). Cost: solve PDE on whole grid (use \textbf{narrow band} or \textbf{fast marching} for speedup). Re-initialize $\phi$ as signed distance periodically.

\textbf{Q13. U-Net — why standard for medical seg?}\\
\textbf{U-Net} (Ronneberger 2015) is a fully-convolutional \textbf{encoder-decoder} with \textbf{skip connections} between symmetric levels. Encoder = contracting path (conv-conv-pool, channels double); decoder = expanding path (up-conv, concat skip, conv-conv). Skips inject \textbf{high-resolution detail} into deep features → sharp boundaries. Reasons it dominates medical: (1) works with \textbf{small datasets} (data augmentation + skips); (2) end-to-end pixel labels; (3) easily extended to 3D (3D U-Net, V-Net); (4) modular — backbone can be swapped (ResNet, EfficientNet, ViT). Variants: Attention U-Net, U-Net++, TransU-Net. The original BraTS 2015 paper already showed U-Net winning.

\textbf{Q14. V-Net vs U-Net.}\\
\textbf{V-Net} (Milletari 2016) is the \textbf{3D} counterpart of U-Net for volumetric data (MRI, CT). Differences: (1) \textbf{3D convolutions} ($k\times k\times k$) instead of 2D; (2) \textbf{residual connections} inside each block (V-Net came before ResU-Net was popular); (3) introduced the \textbf{Dice loss} as primary objective:
\[\mathcal{L}_{\text{Dice}}=1-\frac{2\sum p_i g_i}{\sum p_i^2+\sum g_i^2},\]
which directly handles \textbf{class imbalance} (small foreground); (4) downsampling via \textbf{strided conv} not max-pool. V-Net was originally validated on \textbf{prostate MRI}. Today nnU-Net unifies 2D/3D U-Net/V-Net.

\textbf{Q15. nnU-Net — "no-new-Net".}\\
\textbf{nnU-Net} (Isensee, 2021, Nature Methods) is a \textbf{self-configuring framework}, not a new architecture. The thesis: a properly configured plain U-Net beats fancy custom networks on most medical tasks. Given a dataset, nnU-Net automatically infers \textbf{patch size, batch size, spacing, normalization, network depth, augmentation, training schedule} from a "data fingerprint" (median shape, spacing, modality). It produces 2D / 3D-full / 3D-cascade variants and \textbf{ensembles} the best. Won 19+ segmentation challenges (BraTS, KiTS, AMOS). Lesson: \textbf{good engineering > fancy architecture}. It is the modern baseline you must beat to publish.

\textbf{Q16. Classical vs deep learning segmentation.}\\
\textbf{Classical} (Otsu, region growing, watershed, snakes, level sets, atlas-based, Random Forest+features): interpretable, no training data, fast, parameter-tunable; but brittle, requires hand-crafted features, fails on heterogeneous tumors. \textbf{Deep learning} (U-Net, V-Net, nnU-Net, SAM-Med): end-to-end, state-of-the-art Dice, generalizes within distribution; but needs annotated data (expensive in medicine), black-box, vulnerable to \textbf{domain shift}. In practice: classical for QC / coarse ROI / preprocessing (skull strip), DL for the main segmentation. Hybrid pipelines (DL-init + level-set refine) combine both.

\textbf{Q17. SAM and SAM-Med.}\\
\textbf{SAM} (Segment Anything Model, Meta 2023) is a \textbf{foundation model} for segmentation trained on 1B masks (SA-1B). Architecture: ViT-H image encoder + prompt encoder (point/box/mask) + lightweight mask decoder. Promptable, zero-shot on natural images. \textbf{Limitation on medical}: trained on natural images → poor on MRI/CT/path. \textbf{SAM-Med2D / SAM-Med3D / MedSAM} fine-tune SAM on medical corpora ($\approx$1.5M medical masks). They enable \textbf{interactive segmentation} (click a tumor, get a mask), promising for clinical workflows where fully-automatic is risky. Still struggles on small lesions and cross-modality generalization.

\textbf{Q18. Dice and Jaccard / IoU.}\\
\textbf{Dice} (Sørensen-Dice): $D=\dfrac{2|A\cap B|}{|A|+|B|}=\dfrac{2\,TP}{2\,TP+FP+FN}.$ \textbf{Jaccard / IoU}: $J=\dfrac{|A\cap B|}{|A\cup B|}=\dfrac{TP}{TP+FP+FN}.$ Both range $[0,1]$, 1=perfect. Relation: $D=\dfrac{2J}{1+J}$ and $J=\dfrac{D}{2-D}$; Dice is always $\geq$ Jaccard. Dice is more forgiving (penalizes errors less). Both are \textbf{volume-overlap} metrics — they ignore boundary smoothness. In medicine Dice is the de-facto reporting metric (BraTS, MSD).

\textbf{Q19. Hausdorff95.}\\
\textbf{Hausdorff distance} $H(A,B)=\max\{\sup_{a\in A}\inf_{b\in B}\|a-b\|,\;\sup_{b\in B}\inf_{a\in A}\|a-b\|\}$ measures the \textbf{worst-case boundary error}. It is extremely sensitive to outliers (one stray voxel ruins it). \textbf{HD95} = 95th percentile of the same distances → robust to outliers while still capturing boundary quality. Preferred over Dice when \textbf{boundary accuracy matters}: radiotherapy contours, surgical planning, cortex/vessel segmentation. Reported in millimeters (so spacing matters). BraTS uses HD95 alongside Dice.

\textbf{Q20. ASSD.}\\
\textbf{Average Symmetric Surface Distance}:
\[\text{ASSD}(A,B)=\frac{1}{|\partial A|+|\partial B|}\Big(\sum_{a\in\partial A}d(a,\partial B)+\sum_{b\in\partial B}d(b,\partial A)\Big),\]
the mean distance between surface points of prediction and ground truth, in mm. Symmetric (averaged both ways), unlike one-sided versions. More stable than HD95 because it averages over the whole surface rather than taking a high quantile. Used together with Dice and HD95 in BraTS, ACDC, MSD. Lower = better; sub-millimeter ASSD is excellent for organs, 1--3 mm typical for tumors.

\textbf{Q21. High Dice but poor HD95/ASSD — why?}\\
The metrics measure different things. \textbf{Dice} is a global \textbf{volume overlap}; if 95\% of the voxels are correct Dice is high. \textbf{HD95/ASSD} are \textbf{boundary} metrics. Failure modes: (1) \textbf{Spurious far false positives} — a small predicted blob far from the lesion barely changes Dice but spikes HD95. (2) \textbf{Holes inside} the predicted mask — Dice drops slightly, surface distances can grow if topology is wrong. (3) \textbf{Jagged boundaries} — Dice insensitive, ASSD increases. Paper 1 (Multi-scale 3D-GAN) shows this paradox on \textbf{ET}: Dice up, HD95/ASSD worse → the GAN improves volume but creates boundary noise.

\textbf{Q22. The BraTS challenge.}\\
\textbf{BraTS} (Brain Tumor Segmentation) is a yearly \textbf{MICCAI challenge} (since 2012) on multi-institutional pre-operative MRI of glioma patients, run by U.~Penn (CBICA). It distributes hundreds to thousands of pre-processed (co-registered, skull-stripped, 1\,mm isotropic, $240\!\times\!240\!\times\!155$) multi-modal scans with expert tumor labels. Tasks expanded: segmentation (core), survival prediction, uncertainty, missing-modality (BraTS-2023 has Africa, pediatric, meningioma sub-challenges). It is the \textbf{de facto benchmark} for medical 3D segmentation; nnU-Net, Vox2Vox, Swin-UNetR all compete here.

\textbf{Q23. Four MRI modalities in BraTS.}\\
(1) \textbf{T1} — anatomy, fat bright; (2) \textbf{T1-Gd / T1ce} — post-gadolinium contrast, \textbf{enhancing tumor (ET)} and blood-brain-barrier breakdown light up; (3) \textbf{T2} — water/edema bright; (4) \textbf{FLAIR} — T2 with CSF suppression, \textbf{whole tumor (WT) edema} clearly visible. Multi-modality is essential because each tumor compartment is best seen in a different sequence: \textbf{ET}$\to$T1ce, \textbf{edema/WT}$\to$FLAIR, \textbf{necrosis}$\to$T1, \textbf{tumor core}$\to$T1ce+T2. Networks fuse them as 4-channel input.

\textbf{Q24. WT, TC, ET sub-regions.}\\
BraTS labels are nested: \textbf{label 1} necrotic core, \textbf{label 2} edema, \textbf{label 4} enhancing. Evaluation regions: \textbf{WT (Whole Tumor)} = labels $\{1,2,4\}$ — everything abnormal, primarily on FLAIR. \textbf{TC (Tumor Core)} = $\{1,4\}$ — solid + necrotic, on T1ce. \textbf{ET (Enhancing Tumor)} = $\{4\}$ — actively growing rim on T1ce. ET is the hardest (smallest, irregular), WT the easiest. Reported metrics are Dice/HD95 per region, then averaged. Surgical and radiotherapy planning rely on these sub-regions.

\textbf{Q25. Standard pipeline for brain tumor segmentation.}\\
(1) \textbf{Acquisition}: 4 modalities (T1, T1ce, T2, FLAIR). (2) \textbf{Co-register} to T1 reference (rigid). (3) \textbf{Resample} to 1\,mm isotropic. (4) \textbf{N4 bias correction} on each modality. (5) \textbf{Skull-strip} (HD-BET / ROBEX). (6) \textbf{Intensity normalize} (z-score within brain mask). (7) Stack 4 channels, feed into \textbf{3D U-Net / nnU-Net}. (8) \textbf{Post-process}: connected-component, remove small islands, threshold ET if too few voxels (BraTS trick: zero out ET if $<$small count). (9) Compute Dice/HD95 per WT/TC/ET. (10) Visualize overlays for radiologist QC.

\textbf{Q26. Skull stripping.}\\
\textbf{Skull stripping} (brain extraction) removes non-brain tissue (skull, scalp, eyes, neck) keeping only brain parenchyma. Needed because (a) non-brain has high T1 fat that confuses normalization, (b) reduces ROI for the network. Classical: \textbf{BET} (FSL, deformable surface), \textbf{BSE} (BrainSuite, edge+morphology), \textbf{ROBEX} (random forest + active shape), \textbf{3dSkullStrip} (AFNI). Deep: \textbf{HD-BET} (DKFZ, U-Net trained on multi-modal MRI) — current SOTA, robust to pathologies. BraTS scans come pre-stripped using HD-BET.

\textbf{Q27. End-to-end tumor detection in MRI.}\\
(1) Patient gets multi-modal MRI (T1, T1ce, T2, FLAIR). (2) DICOM $\to$ NIfTI conversion (\texttt{dcm2niix}). (3) Preprocessing: \textbf{co-register} all modalities, \textbf{N4 bias correction}, \textbf{skull-strip} (HD-BET), \textbf{intensity normalize} (z-score). (4) Run \textbf{3D nnU-Net} (or Multi-scale 3D-GAN) trained on BraTS, output 4-class voxel-wise probability. (5) Argmax + connected-component cleanup. (6) Compute volumes for WT/TC/ET, generate \textbf{3D mesh} (marching cubes), overlay on MRI in viewer (3D Slicer). (7) Radiologist reviews; \textbf{XAI heatmap} (Grad-CAM) shows where the network looked. (8) Report exported as DICOM-SR back to PACS for clinical use.

\textbf{Q28. Signals detectable in MRI.}\\
MRI distinguishes tissues by \textbf{T1, T2, proton density, diffusion, perfusion}. Detectable abnormalities: (1) \textbf{Tumor} — mass effect, contrast enhancement on T1ce; (2) \textbf{Edema} — bright on T2/FLAIR around lesions; (3) \textbf{Hemorrhage} — variable depending on age (acute = T2 dark, chronic = T1 bright then T2* bloom); (4) \textbf{Infarct (stroke)} — bright on \textbf{DWI}, dark on ADC, acutely; (5) \textbf{Demyelination} (MS) — periventricular FLAIR plaques; (6) \textbf{Abscess} — ring-enhancing with restricted diffusion inside; (7) \textbf{Microbleeds} — \textbf{SWI} blooming; (8) \textbf{Functional activity} — BOLD fMRI; (9) \textbf{Tracts} — DTI tractography.

\textbf{Q29. Paper 1 summary (Multi-scale 3D-GAN, BraTS 2021).}\\
\textbf{Problem}: brain-tumor segmentation on BraTS 2021. \textbf{Architecture}: Generator = \textbf{3D Residual U-Net} + a novel \textbf{Bridge module} that aggregates multi-scale features before the Discriminator; Discriminator = standard 3D CNN. \textbf{Algorithm 1}: dynamically adjusts the number of \textbf{D-epochs} per G-epoch using the \textbf{linear-regression slope} of adversarial loss → keeps G/D balanced. \textbf{Loss}: \textbf{Dice + MSE adversarial}. \textbf{Results}: Dice ET 82.3, WT 88.6, TC 83.1, beating U-Net, V-Net, VoxelGAN, Vox2Vox baselines. \textbf{Contribution}: GAN with multi-scale Bridge + adaptive D-training.

\textbf{Q30. Bridge module — why useful?}\\
The \textbf{Bridge} sits between Generator's multi-scale feature pyramid and the Discriminator. It \textbf{aggregates features from multiple resolutions} (deep semantic + shallow texture) before D evaluates real-vs-fake. Benefit: D no longer judges only the final mask but also \textbf{intermediate scales}, forcing G to produce masks consistent across scales — helpful for tumors of varying size. It is conceptually similar to \textbf{PatchGAN with FPN}. Weakness in the paper: the design lacks a dedicated figure and there is \textbf{no ablation} isolating its contribution from the GAN baseline.

\textbf{Q31. Algorithm 1 — dynamic D-epoch adjustment.}\\
Classical GAN training is unstable: if D becomes too strong, gradients to G vanish; too weak, G diverges. Algorithm 1 monitors the \textbf{adversarial loss} of D over a sliding window of $l$ recent epochs and computes the \textbf{linear-regression slope} $k$. If the slope is too negative (D rapidly winning, loss dropping fast) → reduce D-epochs. If slope too positive or near zero (D struggling) → increase D-epochs. A threshold $\beta$ separates cases. Net effect: automatic G--D balance without manual scheduling. Paper provides no theoretical justification for the linear-regression choice and no sensitivity study on $k,l,\beta$.

\textbf{Q32. Loss — Dice + MSE, why combine?}\\
The total loss is $\mathcal{L}=\mathcal{L}_{\text{Dice}}(G(x),y)+\lambda\,\mathcal{L}_{\text{MSE-adv}}(D(G(x)),1).$ \textbf{Dice loss} directly optimizes overlap and handles \textbf{class imbalance} (tumor $\ll$ background). \textbf{MSE adversarial} (LSGAN-style) replaces the original cross-entropy GAN loss and gives \textbf{smoother gradients}, fewer vanishing-gradient problems, more stable training. Combining: Dice ensures voxel correctness, adversarial term forces \textbf{globally realistic shapes} (label distribution similar to real masks). Trade-off: GAN may improve Dice but degrade boundary metrics (the ET paradox).

\textbf{Q33. Weaknesses of Paper 1 \& improvements.}\\
(1) Evaluation only on a \textbf{validation split} of training data, no submission to the BraTS \textbf{online server} → not directly comparable to leaderboard. (2) \textbf{No cross-validation} (no mean$\pm$std). (3) \textbf{No ablation} for Bridge or Algorithm 1 individually. (4) Unexplained \textbf{HD95/ASSD paradox} on ET. (5) Hyperparameters $k,l,\beta$ not investigated. (6) No model size, FLOPs, training/inference time. (7) No statistical significance test. (8) No code release. \textbf{Fixes}: 5-fold CV with std, ablation grid, BraTS-server submission, hyperparameter sensitivity, code release on GitHub, parameter / FLOPs table, paired Wilcoxon test.

\textbf{Q34. HD95/ASSD paradox on ET — why an issue?}\\
Multi-scale 3D-GAN reports \textbf{higher Dice} than the ResU-Net (no-GAN) variant on ET, yet \textbf{worse HD95} (9.6mm vs 6.4mm) and \textbf{worse ASSD} (3.37 vs 2.35). This means the GAN \textbf{adds correctly-labeled voxels} (raising overlap) but also produces \textbf{boundary noise / spurious far blobs} (worsening surface distance). For \textbf{enhancing tumor}, where boundaries directly inform \textbf{neurosurgical resection margins} and \textbf{radiotherapy GTV}, boundary errors are clinically dangerous. The paper should have analyzed this — perhaps the MSE adversarial term over-fits to global mass and hallucinates rim voxels. Mitigation: post-process with connected-component, add boundary-aware loss (Boundary loss, HD loss).

\textbf{Q35. Paper 2 summary (ANN$\to$DCNN review, 2024).}\\
\textbf{Thesis}: a "paradigm shift" from \textbf{classic ANNs (with hand-crafted features)} to \textbf{deep CNNs (end-to-end feature learning)} in medical image processing. Key arguments: (i) old pipeline = manual feature extraction (texture, shape, GLCM) + ANN classifier; (ii) new pipeline = raw images $\to$ DCNN that \textbf{learns} features. (iii) Two summary tables compare ANN vs DCNN studies; the "feature extraction" column disappears for DCNN. (iv) Acknowledges DCNN limitations: data-hungry, overfitting, black-box. The paper is mostly \textbf{introductory / pedagogical}.

\textbf{Q36. Strength of Paper 2's Sobel kernel example.}\\
The paper walks through a \textbf{concrete numerical convolution} between a small image patch and a \textbf{Sobel} edge kernel, applied to a hand X-ray. It is computed step by step on actual values, then visualized as the produced edge map. This is excellent \textbf{didactically}: a beginner sees exactly \textbf{what convolution does}, why a CNN's first layers learn edge-like filters, and how feature extraction emerges automatically. It bridges the abstract ("CNN learns features") with the concrete (matrix math), a common gap in survey papers.

\textbf{Q37. Weaknesses of Paper 2.}\\
(1) Published 2024 but \textbf{ignores Vision Transformers, SAM, Med-PaLM, LLaVA-Med} — the actual current paradigm shift. (2) Focuses on \textbf{classification}, almost no \textbf{segmentation} (no U-Net / nnU-Net) and no key datasets (BraTS, LIDC, CheXpert). (3) English/Western-only, \textbf{no low-resource languages}, no Vietnamese-medical-NLP discussion, no medical knowledge graph. (4) No \textbf{XAI} or \textbf{uncertainty quantification} (Grad-CAM, conformal prediction, Dempster-Shafer) despite EU AI Act high-risk classification. (5) No \textbf{PRISMA} review methodology, no inclusion/exclusion criteria, no positioning vs Litjens 2017 / Esteva 2019. (6) Long historical 1980s--90s section, short on current CNN mechanisms; AI/ML/DL definition boxes too basic for a Q1 journal.

\textbf{Q38. Why ignoring ViT and SAM in 2024 is fatal.}\\
By 2024 the field had clearly moved beyond plain CNNs. \textbf{Vision Transformers} (ViT, Swin, MaxViT) and hybrid CNN-Transformers (TransU-Net, Swin-UNetR) dominate medical leaderboards. \textbf{Foundation models} (SAM, MedSAM, Med-PaLM-2, LLaVA-Med) are reshaping how we think about training: pretrain-once, prompt-many. Calling ANN$\to$DCNN the "paradigm shift" of 2024 is anachronistic — the actual shift is \textbf{DCNN $\to$ Transformers / foundation models}. A review missing this is \textbf{out of date on arrival} and gives readers a wrong mental model of where to invest research effort.

\textbf{Q39. Missing: low-resource languages \& XAI.}\\
The paper gives \textbf{zero discussion} of medical AI in non-English settings — no Vietnamese, no Bahasa, no Hindi, no Arabic. For Vietnamese CDSS this is a critical gap: lack of labeled corpora, code-mixed Vietnamese--English--Latin terminology, no Vietnamese medical KG, and ICD-10 terms often loaned. Equally absent: \textbf{XAI} — Grad-CAM, integrated gradients, SHAP for radiomics, conformal prediction for trustworthy intervals. The \textbf{EU AI Act} classifies CDSS as \textbf{high-risk}, mandating explainability and uncertainty. A 2024 medical-AI review without these is incomplete from both technical and regulatory angles.

\textbf{Q40. Overall takeaway from both papers.}\\
Paper 1 (3D-GAN) shows the \textbf{technical-research} pattern: novel idea (Bridge + adaptive D-epochs), competitive Dice, but \textbf{rigor gaps} (no CV, no ablation, no server-validation, paradox unexplained). Paper 2 (ANN$\to$DCNN) shows the \textbf{survey} pitfall: pedagogically good intro, but \textbf{misses the current frontier} (ViT, foundation models, XAI, low-resource). For my own thesis on \textbf{Vietnamese CDSS}, lessons are: (i) always benchmark on official servers with CV, ablations, and code release; (ii) cover modern paradigms (Transformers, SAM-Med); (iii) explicitly address \textbf{low-resource Vietnamese} medical NLP and \textbf{XAI/uncertainty} for high-risk classification under EU AI Act / FDA.

\textbf{Q41. How image processing supports diagnosis.}\\
Image processing turns raw scans into \textbf{quantitative, reproducible biomarkers} that augment human reading. (1) \textbf{Detection} — CADe systems flag suspicious lesions (lung nodules, microcalcifications). (2) \textbf{Segmentation} — gives volume, shape, growth-rate (RECIST). (3) \textbf{Classification} — benign vs malignant, grading. (4) \textbf{Registration / longitudinal} — same patient over time, quantify change. (5) \textbf{Radiomics} — hundreds of texture/shape features feed survival models. (6) \textbf{Decision support} — fuse imaging with EHR, labs, genomics for personalized treatment. The radiologist \textbf{stays in the loop}; AI is a second reader, not a replacement.

\textbf{Q42. End-to-end CDSS pipeline.}\\
(1) \textbf{Data ingestion}: DICOM imaging from PACS, EHR (HL7/FHIR), labs, prescriptions, pathology, genomics. (2) \textbf{Preprocessing}: registration, normalization, NLP on free-text Vietnamese notes. (3) \textbf{Feature extraction}: imaging biomarkers (segmentation $\to$ volumetrics, radiomics) + structured features. (4) \textbf{Model inference}: multi-modal model (e.g., imaging-CNN $+$ tabular MLP $+$ text BERT) outputs risk scores or differentials. (5) \textbf{XAI overlay}: Grad-CAM heatmap, SHAP for tabular, citations for text. (6) \textbf{Uncertainty}: conformal prediction set, calibrated probability. (7) \textbf{Recommendation}: ranked diagnoses + suggested next test, CPG-aligned. (8) \textbf{Clinician UI}: integrated into EMR, one-click accept/override, audit logged. (9) \textbf{Feedback loop}: outcomes captured, used to retrain.

\textbf{Q43. XAI in medical AI — methods.}\\
\textbf{XAI} (eXplainable AI) makes black-box models interpretable; mandatory in medical AI for trust, debugging, regulation. Key methods: (1) \textbf{Grad-CAM} — gradient-weighted class activation maps; visual heatmaps over CT/MRI of "where the model looked". (2) \textbf{Integrated Gradients} / \textbf{SmoothGrad} — pixel attributions. (3) \textbf{SHAP / LIME} — model-agnostic feature importance, common for tabular EHR. (4) \textbf{Attention rollouts} for ViT. (5) \textbf{Conformal prediction} — distribution-free, gives \textbf{prediction sets with calibrated coverage} (e.g., 90\% set "this contains the true diagnosis"). (6) \textbf{Dempster--Shafer} — combine evidence from multi-source. Required by \textbf{EU AI Act} for high-risk CDSS.

\textbf{Q44. Regulatory frameworks for medical AI.}\\
\textbf{USA}: \textbf{FDA 510(k)} (substantial equivalence to predicate device), \textbf{De Novo}, \textbf{PMA} (premarket approval, highest risk), \textbf{Software as a Medical Device (SaMD)} guidance, and the new \textbf{Predetermined Change Control Plan (PCCP)} for AI updates. \textbf{Europe}: \textbf{CE marking under MDR 2017/745} (Medical Device Regulation), with classification I/IIa/IIb/III. \textbf{EU AI Act} (2024) classifies most CDSS as \textbf{high-risk}, demanding risk management, data governance, transparency, human oversight, accuracy, robustness, cybersecurity, post-market monitoring. \textbf{ISO 13485} (QMS), \textbf{IEC 62304} (software lifecycle), \textbf{ISO 14971} (risk). Vietnam relies on Ministry of Health Circular 30/2015 and is gradually adopting ISO frameworks.

\textbf{Q45. Challenges of medical AI in Vietnamese healthcare.}\\
(1) \textbf{Data scarcity}: few public Vietnamese-labeled medical corpora; PACS data is hospital-locked. (2) \textbf{Language}: clinical notes mix \textbf{Vietnamese, English, Latin} terminology (e.g., "U não thể glioblastoma multiforme grade IV"); no widely-adopted Vietnamese medical knowledge graph; ICD-10 mappings inconsistent. (3) \textbf{Infrastructure}: legacy PACS, fragmented EMRs, low GPU availability in provincial hospitals. (4) \textbf{Domain shift}: foreign-trained models (BraTS, ImageNet) underperform on Vietnamese scanners and demographics. (5) \textbf{Regulation}: no clear local AI-medical pathway yet; relying on Circular 30/2015 \& ISO. (6) \textbf{Privacy}: PDPL 2023, but anonymization standards weak. (7) \textbf{Trust}: clinicians need XAI in Vietnamese. (8) \textbf{Workforce}: shortage of medical-informatics specialists. Solutions: federated learning across hospitals, Vietnamese medical KG, local foundation-model fine-tunes (PhoBERT-Med, ViMedSAM).

\switchcolumn

\section*{QA3b. Câu hỏi lý thuyết — Phân vùng, BraTS \& Bài báo}

\textbf{Câu 1. Pipeline đầy đủ từ bệnh nhân đến ảnh lâm sàng.}\\
Chuỗi gồm \textbf{7 bước}: (1) \textbf{Acquisition} — bệnh nhân được chụp CT/MRI/US/PET, thu dữ liệu thô \textbf{k-space} hoặc \textbf{sinogram}. (2) \textbf{Reconstruction} — Fourier nghịch (MRI) hay filtered back-projection (CT) dựng volume. (3) \textbf{Tiền xử lý} — \textbf{N4 bias correction}, khử nhiễu, \textbf{intensity normalization} (Nyul), \textbf{skull-stripping}. (4) \textbf{Registration} — căn chỉnh các modality hoặc atlas về không gian bệnh nhân. (5) \textbf{Segmentation} — tách cơ quan / tổn thương (U-Net, nnU-Net). (6) \textbf{Định lượng} — thể tích, hình dáng, radiomics. (7) \textbf{Báo cáo / CDSS} — visualization, DICOM-SR, hỗ trợ quyết định cho bác sĩ. Mọi bước đều propagate lỗi nên QC ở mỗi tầng là bắt buộc.

\textbf{Câu 2. N4 bias correction — vì sao MRI cần?}\\
\textbf{Bias field} là độ không đồng nhất cường độ \textbf{nhân tính (multiplicative)} tần số thấp do \textbf{B1 RF coil} không đồng đều trên MRI; cùng một mô (ví dụ chất trắng) lại sáng bên này, tối bên kia. \textbf{N4ITK} (Tustison 2010) cải tiến từ \textbf{N3}, mô hình hóa log-bias bằng mặt \textbf{B-spline} và lặp deconvolution để histogram sắc nét hơn. Không có N4, segmentation theo cường độ (Otsu, k-means, kể cả CNN) sẽ thất bại vì cùng class lại có cường độ khác nhau. Trong BraTS, N4 là bước bắt buộc. CT không cần (đơn vị Hounsfield đã được hiệu chuẩn).

\textbf{Câu 3. Registration — rigid, affine, B-spline, Demons, LDDMM.}\\
\textbf{Registration} căn ảnh $I$ về ảnh tham chiếu $J$ bằng cách tìm transform $T$ tối thiểu hóa khoảng cách. \textbf{Rigid} (6 DoF: 3 quay + 3 tịnh tiến) — cùng bệnh nhân, cùng modality, theo dõi theo thời gian. \textbf{Affine} (12 DoF) — thêm scaling + shear, dùng cho căn chỉnh thô liên-bệnh nhân. \textbf{B-spline FFD} — non-rigid, biến dạng theo lưới control point, mượt nhưng không đảm bảo invertible. \textbf{Demons} (Thirion) — kiểu optical-flow, nhanh, regularize bằng làm mượt Gaussian. \textbf{LDDMM} — geodesic trên không gian diffeomorphism, đảm bảo \textbf{bảo toàn topology} (không bị gấp nếp), dùng trong computational anatomy.

\textbf{Câu 4. Metric tương đồng trong registration.}\\
Tùy modality. \textbf{SSD} $\sum (I-J)^2$ — chỉ dùng được cùng modality, cùng thang. \textbf{NCC} (normalized cross-correlation) — bền với biến đổi tuyến tính, dùng MRI cùng modality. \textbf{Mutual Information} \(\mathrm{MI}(I,J)=\sum p(i,j)\log\frac{p(i,j)}{p(i)p(j)}\) — chuẩn vàng cho \textbf{multi-modal} (MRI$\leftrightarrow$CT, T1$\leftrightarrow$T2) vì chỉ giả định phụ thuộc thống kê, không cần intensity bằng nhau. \textbf{NMI} (chuẩn hóa MI) giảm bias do overlap. Với deep registration (VoxelMorph) thường dùng \textbf{NCC + smoothness loss} trên trường dịch chuyển.

\textbf{Câu 5. Phương pháp khử nhiễu phổ biến.}\\
\textbf{Gaussian} — đơn giản, làm mờ cạnh. \textbf{Median} — trị nhiễu muối-tiêu, giữ cạnh, dùng cho nhiễu speckle siêu âm. \textbf{Bilateral} — giữ cạnh nhờ kernel kết hợp không gian + cường độ. \textbf{Anisotropic diffusion} (Perona-Malik) — PDE khuếch tán dọc theo cạnh, không qua cạnh. \textbf{Non-Local Means (NLM)} — trung bình các patch tương tự khắp ảnh, là chuẩn vàng tiền-DL. \textbf{BM3D} — block-matching + 3D collaborative filter, gần tối ưu với nhiễu Gaussian. \textbf{DnCNN / Noise2Noise} — deep learning, học residual nhiễu. MRI cần mô hình \textbf{Rician noise}.

\textbf{Câu 6. Nyul normalization — vì sao cần xuyên scanner?}\\
Cường độ MRI \textbf{không được hiệu chuẩn} (khác CT có Hounsfield): cùng não nhưng máy Siemens 3T và Philips 1.5T cho dải cường độ rất khác. Điều này phá vỡ generalization. \textbf{Nyul \& Udupa (2000)} chọn một tập \textbf{landmark} histogram (percentile 1, 10, 20, ..., 99) trên corpus train, tính vị trí trung bình, rồi \textbf{ánh xạ piecewise-linear} mỗi scan mới về các landmark đó. Lựa chọn khác: \textbf{z-score} từng volume, \textbf{white-stripe} (chuẩn theo đỉnh chất trắng), \textbf{histogram matching}. Không có normalization, model train ở bệnh viện A sẽ thất bại ở bệnh viện B — bài toán \textbf{domain shift} kinh điển.

\textbf{Câu 7. Segmentation là gì? Vì sao quan trọng nhất?}\\
\textbf{Segmentation} gán nhãn cho từng voxel: $S:\Omega\to\{0,1,\ldots,K\}$, chia ảnh thành các vùng có ý nghĩa giải phẫu/bệnh lý. Đây là cầu nối từ pixel sang \textbf{biomarker định lượng}: \textbf{thể tích} u, tốc độ phát triển (RECIST), hình dáng, radiomics. Chẩn đoán ("có u không?"), lập kế hoạch xạ trị (vẽ contour), dẫn hướng phẫu thuật, dự báo tiên lượng — đều cần segmentation. Classification chỉ trả lời "có/không"; segmentation trả lời "ở đâu + lớn cỡ nào". Vì vậy BraTS, LiTS, KiTS đều xoay quanh segmentation.

\textbf{Câu 8. Otsu — công thức.}\\
\textbf{Otsu (1979)} chọn ngưỡng $t^*$ \textbf{tối đa hóa phương sai giữa-lớp} (tương đương tối thiểu phương sai trong-lớp) trên histogram. Với xác suất $\omega_0(t),\omega_1(t)$ và trung bình $\mu_0(t),\mu_1(t)$:
\[t^*=\arg\max_t \;\sigma_B^2(t)=\omega_0(t)\,\omega_1(t)\,[\mu_0(t)-\mu_1(t)]^2.\]
Giả định histogram \textbf{lưỡng phương thức} (foreground + background). Ưu: không tham số, $O(L)$. Nhược: thất bại nếu nhiều mode, nhạy với nhiễu và chiếu sáng không đều (dùng \textbf{Otsu cục bộ / adaptive}). Có biến thể \textbf{multi-Otsu} cho nhiều class.

\textbf{Câu 9. Region growing.}\\
\textbf{Region growing} bắt đầu từ một (hoặc vài) \textbf{pixel mầm} và liên tục hợp nhất các neighbor thoả tiêu chí \textbf{đồng nhất} (ví dụ $|I(p)-\mu_R|<\tau$ hay theo gradient). Đây là phương pháp \textbf{bán tự động} (người dùng chọn seed), trực quan, từng được dùng trong segment gan/phổi trên các tool PACS đời cũ. Nhược: rất nhạy seed, hay bị \textbf{rò rỉ (leakage)} khi biên yếu, không tách được hai cấu trúc dính nhau, ngưỡng $\tau$ khó tinh chỉnh. Biến thể: \textbf{symmetric region growing}, \textbf{statistical region merging}, \textbf{confidence connected} (ITK).

\textbf{Câu 10. Watershed — vì sao over-segment, cách khắc phục.}\\
\textbf{Watershed} coi ảnh gradient là bề mặt địa hình; "ngập nước" từ các cực tiểu cục bộ tạo các bồn chứa, ngăn cách bởi \textbf{watershed line}. Mỗi bồn = một vùng. Vấn đề: \textbf{mỗi cực tiểu cục bộ thành một vùng} → \textbf{over-segmentation} nghiêm trọng trên gradient nhiễu (hàng trăm vùng). Khắc phục: (1) \textbf{Marker-controlled watershed} — chỉ ngập từ marker do người/thuật toán cung cấp; (2) \textbf{H-minima transform} — bỏ các cực tiểu nông; (3) \textbf{Smooth} gradient trước (Gaussian/bilateral); (4) \textbf{Region merging} hậu xử lý. Marker-controlled là biến thể chuẩn trong segmentation tế bào.

\textbf{Câu 11. Snakes — công thức năng lượng.}\\
\textbf{Snake} là đường tham số $C(s)=(x(s),y(s)),\,s\in[0,1]$ tiến hoá để tối thiểu:
\[E_{\text{snake}}=\int_0^1\!\!\big(\underbrace{\alpha|C'(s)|^2+\beta|C''(s)|^2}_{E_{\text{int}}}+\underbrace{E_{\text{ext}}(C(s))}_{-|\nabla I|^2}\big)\,ds.\]
\textbf{Năng lượng nội}: $\alpha$ điều khiển \textbf{độ đàn hồi}, $\beta$ điều khiển \textbf{độ cứng (chống cong)}. \textbf{Năng lượng ngoài} kéo đường về phía cạnh ($|\nabla I|$ lớn). Giải bằng gradient descent qua Euler-Lagrange sai phân hữu hạn. Nhược: cần khởi tạo tốt, không đổi topology, kẹt ở vùng lõm. \textbf{GVF snake} mở rộng vùng hấp dẫn.

\textbf{Câu 12. Level set — PDE tiến hoá.}\\
\textbf{Level set} (Osher-Sethian 1988) nhúng đường viền vào mức 0 của hàm cao chiều $\phi(\mathbf{x},t)$: $C(t)=\{\mathbf{x}:\phi(\mathbf{x},t)=0\}$. PDE tiến hoá:
\[\frac{\partial\phi}{\partial t}+F\,|\nabla\phi|=0,\]
với tốc độ $F$ phụ thuộc curvature $\kappa=\nabla\!\cdot\!(\nabla\phi/|\nabla\phi|)$ và đặc trưng ảnh. \textbf{Chan-Vese} dùng thống kê vùng (mean trong/ngoài) thay cạnh. Lợi thế lớn so với snake: \textbf{tự động đổi topology} (tách/gộp). Chi phí: giải PDE trên cả lưới (dùng \textbf{narrow band} hoặc \textbf{fast marching} để tăng tốc). Cần re-init $\phi$ thành signed distance định kỳ.

\textbf{Câu 13. U-Net — vì sao chuẩn cho y tế?}\\
\textbf{U-Net} (Ronneberger 2015) là kiến trúc fully-convolutional dạng \textbf{encoder-decoder} có \textbf{skip connection} đối xứng. Encoder = đường thu hẹp (conv-conv-pool, channel gấp đôi); decoder = đường mở rộng (up-conv, concat skip, conv-conv). Skip "tiêm" \textbf{chi tiết độ phân giải cao} vào feature sâu → biên sắc nét. Lý do thống trị y tế: (1) hiệu quả với \textbf{dataset nhỏ} (augmentation + skip); (2) end-to-end, ra label per-pixel; (3) dễ mở rộng 3D (3D U-Net, V-Net); (4) backbone modular (ResNet, EfficientNet, ViT). Biến thể: Attention U-Net, U-Net++, TransU-Net. BraTS 2015 đã thấy U-Net thắng.

\textbf{Câu 14. V-Net khác U-Net.}\\
\textbf{V-Net} (Milletari 2016) là phiên bản \textbf{3D} của U-Net cho dữ liệu volumetric (MRI, CT). Khác biệt: (1) \textbf{conv 3D} ($k\times k\times k$) thay vì 2D; (2) \textbf{residual connection} bên trong mỗi block (V-Net ra trước ResU-Net); (3) đề xuất \textbf{Dice loss} làm objective chính:
\[\mathcal{L}_{\text{Dice}}=1-\frac{2\sum p_i g_i}{\sum p_i^2+\sum g_i^2},\]
xử lý trực tiếp \textbf{class imbalance} (foreground rất nhỏ); (4) downsampling bằng \textbf{strided conv} thay max-pool. V-Net được kiểm chứng đầu tiên trên \textbf{MRI tuyến tiền liệt}. Ngày nay nnU-Net thống nhất 2D/3D U-Net/V-Net.

\textbf{Câu 15. nnU-Net — "no-new-Net".}\\
\textbf{nnU-Net} (Isensee, Nature Methods 2021) là \textbf{framework tự cấu hình}, không phải kiến trúc mới. Luận điểm: U-Net cấu hình đúng cách thắng các kiến trúc cầu kỳ trên hầu hết tác vụ y tế. Cho một dataset, nnU-Net tự suy ra \textbf{patch size, batch size, spacing, normalization, depth, augmentation, training schedule} từ "data fingerprint" (median shape, spacing, modality). Nó tạo các biến thể 2D / 3D-full / 3D-cascade rồi \textbf{ensemble} cái tốt nhất. Thắng 19+ challenge (BraTS, KiTS, AMOS). Bài học: \textbf{kỹ thuật tốt > kiến trúc fancy}. Đây là baseline bạn phải đánh bại để công bố.

\textbf{Câu 16. Cổ điển vs deep learning.}\\
\textbf{Cổ điển} (Otsu, region growing, watershed, snake, level set, atlas, Random Forest+feature): dễ giải thích, không cần data, nhanh, có thể tinh chỉnh; nhưng dễ vỡ, cần feature thủ công, thất bại với u không đồng nhất. \textbf{Deep learning} (U-Net, V-Net, nnU-Net, SAM-Med): end-to-end, Dice cao nhất, generalize tốt trong phân phối; nhưng cần dữ liệu nhãn (rất đắt trong y tế), black-box, dễ tổn thương \textbf{domain shift}. Thực tế: cổ điển dùng cho QC / khoanh ROI thô / preprocessing (skull strip), DL cho segment chính. Pipeline lai (DL khởi tạo + level-set tinh chỉnh) kết hợp cả hai.

\textbf{Câu 17. SAM và SAM-Med.}\\
\textbf{SAM} (Segment Anything, Meta 2023) là \textbf{foundation model} cho segmentation, train trên 1B mask (SA-1B). Kiến trúc: ViT-H image encoder + prompt encoder (point/box/mask) + mask decoder nhẹ. Promptable, zero-shot trên ảnh tự nhiên. \textbf{Hạn chế y tế}: train trên ảnh tự nhiên → kém trên MRI/CT/giải phẫu bệnh. \textbf{SAM-Med2D / SAM-Med3D / MedSAM} fine-tune SAM trên corpus y tế (\(\approx\)1.5M mask). Cho phép \textbf{segmentation tương tác} (click vào u $\to$ ra mask), hứa hẹn cho workflow lâm sàng nơi tự động hoàn toàn còn rủi ro. Vẫn yếu với tổn thương nhỏ và generalize liên modality.

\textbf{Câu 18. Dice và Jaccard / IoU.}\\
\textbf{Dice}: $D=\dfrac{2|A\cap B|}{|A|+|B|}=\dfrac{2\,TP}{2\,TP+FP+FN}.$ \textbf{Jaccard / IoU}: $J=\dfrac{|A\cap B|}{|A\cup B|}=\dfrac{TP}{TP+FP+FN}.$ Cả hai $\in[0,1]$, 1=hoàn hảo. Quan hệ: $D=\dfrac{2J}{1+J}$, $J=\dfrac{D}{2-D}$; Dice luôn $\geq$ Jaccard. Dice "dễ tính" hơn (phạt lỗi nhẹ hơn). Cả hai là \textbf{overlap thể tích} — không quan tâm độ mượt biên. Trong y tế Dice là metric báo cáo chuẩn (BraTS, MSD).

\textbf{Câu 19. Hausdorff95.}\\
\textbf{Khoảng cách Hausdorff} $H(A,B)=\max\{\sup_{a\in A}\inf_{b\in B}\|a-b\|,\;\sup_{b\in B}\inf_{a\in A}\|a-b\|\}$ đo \textbf{lỗi biên tệ nhất}. Cực kỳ nhạy với outlier (một voxel lạc cũng làm hỏng). \textbf{HD95} = phân vị 95\% của cùng các khoảng cách → bền với outlier mà vẫn phản ánh chất lượng biên. Ưu tiên hơn Dice khi \textbf{độ chính xác biên} quan trọng: vẽ contour xạ trị, lập kế hoạch phẫu thuật, segment vỏ não/mạch máu. Đơn vị mm (nên spacing quan trọng). BraTS dùng HD95 cùng Dice.

\textbf{Câu 20. ASSD.}\\
\textbf{Average Symmetric Surface Distance}:
\[\text{ASSD}(A,B)=\frac{1}{|\partial A|+|\partial B|}\Big(\sum_{a\in\partial A}d(a,\partial B)+\sum_{b\in\partial B}d(b,\partial A)\Big),\]
khoảng cách trung bình giữa điểm bề mặt dự đoán và ground truth, đơn vị mm. Đối xứng (trung bình hai chiều). Ổn định hơn HD95 vì lấy trung bình toàn bề mặt, không phải phân vị cao. Dùng cùng Dice và HD95 trong BraTS, ACDC, MSD. Nhỏ = tốt; ASSD dưới 1mm là xuất sắc với cơ quan, 1--3mm là điển hình với u.

\textbf{Câu 21. Dice cao nhưng HD95/ASSD tệ — vì sao?}\\
Hai metric đo hai thứ khác nhau. \textbf{Dice} là \textbf{overlap thể tích} toàn cục; 95\% voxel đúng là Dice cao. \textbf{HD95/ASSD} là metric \textbf{biên}. Các kiểu lỗi: (1) \textbf{FP nhỏ ở xa} — không ảnh hưởng Dice nhưng đẩy HD95 lên. (2) \textbf{Lỗ rỗng bên trong} mask — Dice giảm nhẹ, surface distance tăng nếu topology sai. (3) \textbf{Biên răng cưa} — Dice không nhạy, ASSD tăng. Bài 1 (Multi-scale 3D-GAN) cho thấy nghịch lý này trên \textbf{ET}: Dice tăng nhưng HD95/ASSD tệ hơn → GAN cải thiện thể tích nhưng tạo nhiễu biên.

\textbf{Câu 22. BraTS challenge.}\\
\textbf{BraTS} (Brain Tumor Segmentation) là \textbf{thử thách MICCAI} hằng năm (từ 2012) trên ảnh MRI tiền phẫu đa-trung-tâm của bệnh nhân glioma, do U.~Penn (CBICA) tổ chức. Phát hành hàng trăm đến hàng ngàn ca đã được tiền xử lý (đồng đăng ký, skull-strip, isotropic 1\,mm, $240\!\times\!240\!\times\!155$) đa modality kèm nhãn chuyên gia. Mở rộng dần: segmentation (lõi), dự báo sống, uncertainty, missing-modality (BraTS-2023 có sub-challenge châu Phi, nhi khoa, meningioma). Đây là \textbf{benchmark de facto} cho segmentation 3D y tế; nnU-Net, Vox2Vox, Swin-UNetR đều cạnh tranh ở đây.

\textbf{Câu 23. Bốn modality MRI trong BraTS.}\\
(1) \textbf{T1} — giải phẫu, mỡ sáng; (2) \textbf{T1-Gd / T1ce} — sau tiêm gadolinium, \textbf{enhancing tumor (ET)} và phá vỡ hàng rào máu não bắt thuốc; (3) \textbf{T2} — nước/phù sáng; (4) \textbf{FLAIR} — T2 ức chế CSF, \textbf{whole tumor (WT) / phù} hiện rõ. Đa modality cần thiết vì mỗi vùng u rõ nhất ở một sequence: \textbf{ET}$\to$T1ce, \textbf{phù/WT}$\to$FLAIR, \textbf{hoại tử}$\to$T1, \textbf{lõi u}$\to$T1ce+T2. Mạng feed 4 channel.

\textbf{Câu 24. WT, TC, ET.}\\
Nhãn BraTS lồng nhau: \textbf{label 1} hoại tử, \textbf{label 2} phù, \textbf{label 4} enhancing. Vùng đánh giá: \textbf{WT (Whole Tumor)} = $\{1,2,4\}$ — toàn bộ bất thường, chủ yếu trên FLAIR. \textbf{TC (Tumor Core)} = $\{1,4\}$ — phần đặc + hoại tử, trên T1ce. \textbf{ET (Enhancing Tumor)} = $\{4\}$ — viền đang phát triển trên T1ce. ET khó nhất (nhỏ, không đều), WT dễ nhất. Báo cáo Dice/HD95 từng vùng rồi trung bình. Phẫu thuật và xạ trị dựa trên các vùng này.

\textbf{Câu 25. Pipeline chuẩn segment u não.}\\
(1) \textbf{Chụp}: 4 modality (T1, T1ce, T2, FLAIR). (2) \textbf{Đồng đăng ký} về T1 (rigid). (3) \textbf{Resample} về isotropic 1\,mm. (4) \textbf{N4 bias correction} từng modality. (5) \textbf{Skull-strip} (HD-BET / ROBEX). (6) \textbf{Chuẩn hóa cường độ} (z-score trong mask não). (7) Stack 4 channel, đẩy vào \textbf{3D U-Net / nnU-Net}. (8) \textbf{Hậu xử lý}: connected-component, bỏ các đảo nhỏ, threshold ET nếu quá ít voxel (mẹo BraTS: zero ET nếu nhỏ). (9) Tính Dice/HD95 cho WT/TC/ET. (10) Visualize overlay cho bác sĩ X quang QC.

\textbf{Câu 26. Skull stripping.}\\
\textbf{Skull stripping} (brain extraction) loại bỏ mô không phải não (sọ, da đầu, mắt, cổ), chỉ giữ nhu mô não. Cần thiết vì (a) mỡ ngoài não T1 cao gây nhiễu normalization, (b) thu hẹp ROI cho mạng. Cổ điển: \textbf{BET} (FSL, mặt biến dạng), \textbf{BSE} (BrainSuite, cạnh + morphology), \textbf{ROBEX} (random forest + active shape), \textbf{3dSkullStrip} (AFNI). Deep: \textbf{HD-BET} (DKFZ, U-Net trên đa modality MRI) — SOTA hiện nay, bền với bệnh lý. Dữ liệu BraTS đã được skull-strip sẵn bằng HD-BET.

\textbf{Câu 27. Quy trình end-to-end phát hiện u não.}\\
(1) Bệnh nhân chụp MRI đa modality (T1, T1ce, T2, FLAIR). (2) DICOM $\to$ NIfTI (\texttt{dcm2niix}). (3) Tiền xử lý: \textbf{đồng đăng ký} các modality, \textbf{N4}, \textbf{skull-strip} (HD-BET), \textbf{chuẩn hóa cường độ} (z-score). (4) Chạy \textbf{3D nnU-Net} (hoặc Multi-scale 3D-GAN) đã train BraTS, ra xác suất 4 lớp per-voxel. (5) Argmax + clean connected-component. (6) Tính thể tích WT/TC/ET, dựng \textbf{lưới 3D} bằng marching cubes, overlay lên MRI ở viewer (3D Slicer). (7) Bác sĩ duyệt; \textbf{XAI heatmap} (Grad-CAM) cho thấy mạng đã "nhìn" đâu. (8) Báo cáo xuất DICOM-SR trở lại PACS để dùng lâm sàng.

\textbf{Câu 28. Tín hiệu MRI có thể phát hiện.}\\
MRI phân biệt mô bằng \textbf{T1, T2, mật độ proton, khuếch tán, tưới máu}. Các bất thường: (1) \textbf{U} — khối choáng chỗ, bắt thuốc T1ce; (2) \textbf{Phù} — sáng T2/FLAIR quanh tổn thương; (3) \textbf{Xuất huyết} — biến đổi theo tuổi (cấp = T2 tối, mãn = T1 sáng rồi T2* "bloom"); (4) \textbf{Nhồi máu (đột quỵ)} — sáng \textbf{DWI}, tối ADC, ở giai đoạn cấp; (5) \textbf{Mất myelin} (MS) — mảng FLAIR cạnh não thất; (6) \textbf{Áp xe} — viền bắt thuốc với DWI hạn chế bên trong; (7) \textbf{Vi xuất huyết} — \textbf{SWI} bloom; (8) \textbf{Hoạt động chức năng} — fMRI BOLD; (9) \textbf{Sợi trục} — DTI tractography.

\textbf{Câu 29. Tóm tắt Bài 1 (Multi-scale 3D-GAN, BraTS 2021).}\\
\textbf{Bài toán}: segment u não trên BraTS 2021. \textbf{Kiến trúc}: Generator = \textbf{3D Residual U-Net} + module \textbf{Bridge} mới gộp đặc trưng đa-tỷ-lệ trước khi đưa vào Discriminator; D là CNN 3D chuẩn. \textbf{Algorithm 1}: tự điều chỉnh số \textbf{epoch của D} theo \textbf{độ dốc hồi quy tuyến tính} của adversarial loss → giữ G/D cân bằng. \textbf{Loss}: \textbf{Dice + MSE adversarial}. \textbf{Kết quả}: Dice ET 82.3, WT 88.6, TC 83.1, vượt U-Net, V-Net, VoxelGAN, Vox2Vox. \textbf{Đóng góp}: GAN có Bridge đa tỷ lệ + huấn luyện D thích ứng.

\textbf{Câu 30. Module Bridge — vì sao hữu ích?}\\
\textbf{Bridge} đặt giữa kim tự tháp đặc trưng đa tỷ lệ của G và Discriminator. Nó \textbf{gộp đặc trưng nhiều độ phân giải} (ngữ nghĩa sâu + texture nông) trước khi D đánh giá real-vs-fake. Lợi: D không chỉ đánh giá mask cuối mà cả \textbf{tỷ lệ trung gian}, ép G tạo mask nhất quán xuyên tỷ lệ — hữu ích cho u kích thước đa dạng. Về ý tưởng giống \textbf{PatchGAN với FPN}. Yếu điểm bài viết: thiếu hình minh họa riêng và \textbf{không có ablation} tách đóng góp Bridge khỏi GAN baseline.

\textbf{Câu 31. Algorithm 1 — điều chỉnh epoch D động.}\\
GAN train không ổn: D quá mạnh → gradient đến G tiêu biến; D quá yếu → G phân kỳ. Algorithm 1 theo dõi \textbf{adversarial loss} của D trên cửa sổ trượt $l$ epoch gần nhất, tính \textbf{độ dốc hồi quy tuyến tính} $k$. Nếu $k$ âm sâu (D thắng nhanh, loss giảm mạnh) → giảm số epoch D. Nếu $k$ dương hoặc gần 0 (D đang đuối) → tăng epoch D. Ngưỡng $\beta$ phân biệt. Hệ quả: tự động cân bằng G--D không cần lịch tay. Bài thiếu \textbf{biện minh lý thuyết} cho việc dùng hồi quy tuyến tính và không khảo sát độ nhạy $k,l,\beta$.

\textbf{Câu 32. Loss Dice + MSE — vì sao kết hợp?}\\
Tổng loss $\mathcal{L}=\mathcal{L}_{\text{Dice}}(G(x),y)+\lambda\,\mathcal{L}_{\text{MSE-adv}}(D(G(x)),1).$ \textbf{Dice loss} tối ưu trực tiếp overlap và xử lý \textbf{class imbalance} (u $\ll$ background). \textbf{MSE adversarial} (kiểu LSGAN) thay cross-entropy GAN gốc, cho \textbf{gradient mượt hơn}, ít vanishing-gradient, train ổn định. Kết hợp: Dice đảm bảo voxel đúng, term adversarial ép \textbf{hình dạng tổng thể realistic} (phân phối nhãn giống mask thật). Đánh đổi: GAN có thể tăng Dice nhưng làm xấu metric biên (nghịch lý ET).

\textbf{Câu 33. Yếu điểm Bài 1 \& cải thiện.}\\
(1) Đánh giá chỉ trên \textbf{validation split} của tập train, không submit \textbf{server BraTS chính thức} → không so sánh trực tiếp leaderboard. (2) \textbf{Không cross-validation} (không có mean$\pm$std). (3) \textbf{Không ablation} riêng Bridge và Algorithm 1. (4) \textbf{Nghịch lý HD95/ASSD} trên ET không được giải thích. (5) Hyperparameter $k,l,\beta$ chưa khảo sát. (6) Không báo cáo size model, FLOPs, thời gian train/inference. (7) Không kiểm định ý nghĩa thống kê. (8) Không công bố code. \textbf{Sửa}: 5-fold CV với std, lưới ablation, submit server BraTS, khảo sát độ nhạy hyperparameter, code GitHub, bảng tham số / FLOPs, kiểm định Wilcoxon ghép cặp.

\textbf{Câu 34. Nghịch lý HD95/ASSD trên ET — vì sao quan trọng?}\\
Multi-scale 3D-GAN báo \textbf{Dice cao hơn} ResU-Net (không GAN) trên ET, nhưng \textbf{HD95 tệ hơn} (9.6mm vs 6.4mm) và \textbf{ASSD tệ hơn} (3.37 vs 2.35). Nghĩa là GAN \textbf{thêm đúng voxel} (tăng overlap) nhưng cũng tạo \textbf{nhiễu biên / đốm xa giả}. Với \textbf{enhancing tumor}, biên trực tiếp xác định \textbf{ranh giới cắt phẫu thuật} và \textbf{GTV xạ trị} → lỗi biên nguy hiểm về lâm sàng. Bài đáng lẽ phải phân tích — có thể MSE adversarial overfit khối tổng thể và "ảo giác" voxel viền. Giảm thiểu: hậu xử lý connected-component, thêm boundary-aware loss (Boundary loss, HD loss).

\textbf{Câu 35. Tóm tắt Bài 2 (review ANN$\to$DCNN, 2024).}\\
\textbf{Luận điểm}: "paradigm shift" từ \textbf{ANN cổ điển (feature thủ công)} sang \textbf{deep CNN (học feature end-to-end)} trong xử lý ảnh y tế. Lập luận chính: (i) pipeline cũ = trích đặc trưng thủ công (texture, hình dạng, GLCM) + ANN classifier; (ii) pipeline mới = ảnh thô $\to$ DCNN \textbf{tự học} feature; (iii) Hai bảng tổng hợp so sánh nghiên cứu ANN và DCNN; cột "feature extraction" \textbf{biến mất} ở DCNN. (iv) Thừa nhận hạn chế DCNN: cần dữ liệu lớn, dễ overfit, black-box. Bài chủ yếu mang tính \textbf{nhập môn / sư phạm}.

\textbf{Câu 36. Điểm mạnh ví dụ Sobel của Bài 2.}\\
Bài đi qua một phép \textbf{tích chập số cụ thể} giữa patch ảnh nhỏ và kernel \textbf{Sobel}, áp lên ảnh X-quang bàn tay. Tính từng bước trên giá trị thật, sau đó hiển thị kết quả là bản đồ cạnh. Đây là \textbf{rất tốt cho dạy học}: người mới thấy chính xác \textbf{convolution làm gì}, vì sao lớp đầu CNN học các filter giống cạnh, và feature extraction "tự nhiên" xuất hiện ra sao. Cầu nối từ trừu tượng ("CNN học đặc trưng") sang cụ thể (toán ma trận) — khoảng trống thường gặp ở các paper review.

\textbf{Câu 37. Yếu điểm Bài 2.}\\
(1) Xuất bản 2024 nhưng \textbf{bỏ qua Vision Transformer, SAM, Med-PaLM, LLaVA-Med} — đúng paradigm shift hiện nay. (2) Chủ yếu \textbf{classification}, gần như không có \textbf{segmentation} (không U-Net / nnU-Net) và không nhắc dataset chính (BraTS, LIDC, CheXpert). (3) Chỉ nói tiếng Anh / phương Tây, \textbf{không đề cập low-resource language}, không nói NLP y tế tiếng Việt, không có medical knowledge graph. (4) Không \textbf{XAI} hay \textbf{uncertainty} (Grad-CAM, conformal prediction, Dempster-Shafer) dù EU AI Act xếp CDSS là high-risk. (5) Không \textbf{PRISMA}, không tiêu chí lựa chọn, không định vị so với Litjens 2017 / Esteva 2019. (6) Phần lịch sử 1980s--90s dài, phần cơ chế CNN ngắn; box định nghĩa AI/ML/DL quá cơ bản với tạp chí Q1.

\textbf{Câu 38. Vì sao bỏ qua ViT, SAM năm 2024 là sai lầm chí mạng.}\\
Đến 2024 ngành đã rõ ràng vượt khỏi CNN thuần. \textbf{Vision Transformer} (ViT, Swin, MaxViT) và CNN-Transformer lai (TransU-Net, Swin-UNetR) thống trị bảng xếp hạng y tế. \textbf{Foundation model} (SAM, MedSAM, Med-PaLM-2, LLaVA-Med) đang định hình lại cách train: pretrain một lần, prompt nhiều lần. Gọi ANN$\to$DCNN là "paradigm shift" của 2024 là lạc hậu — dịch chuyển thực sự đang là \textbf{DCNN $\to$ Transformer / foundation model}. Một review thiếu phần này là \textbf{lỗi thời ngay khi xuất bản} và đưa ra mô hình tư duy sai cho người đọc về nơi nên đầu tư nghiên cứu.

\textbf{Câu 39. Thiếu sót: low-resource language \& XAI.}\\
Bài viết \textbf{không hề bàn} về AI y tế ngoài bối cảnh tiếng Anh — không Việt, không Bahasa, không Hindi, không Ả-rập. Với CDSS tiếng Việt đây là khoảng trống chí tử: thiếu corpus có nhãn, thuật ngữ pha trộn Việt-Anh-Latin, không có KG y tế tiếng Việt, ICD-10 thường vay mượn. Tương tự thiếu: \textbf{XAI} — Grad-CAM, integrated gradients, SHAP cho radiomics, conformal prediction cho khoảng tin cậy. \textbf{EU AI Act} xếp CDSS là \textbf{high-risk}, bắt buộc explainability và uncertainty. Một review y-tế-AI 2024 thiếu các phần này là không đầy đủ cả về kỹ thuật lẫn pháp lý.

\textbf{Câu 40. Bài học chung từ hai bài.}\\
Bài 1 (3D-GAN) cho thấy mẫu của \textbf{nghiên cứu kỹ thuật}: ý tưởng mới (Bridge + epoch D thích ứng), Dice cạnh tranh, nhưng \textbf{thiếu chặt chẽ} (không CV, không ablation, không server-validate, không giải thích nghịch lý). Bài 2 (ANN$\to$DCNN) cho thấy cạm bẫy của \textbf{survey}: nhập môn tốt, nhưng \textbf{lỡ biên giới hiện tại} (ViT, foundation model, XAI, low-resource). Với luận văn của tôi về \textbf{CDSS tiếng Việt}, bài học: (i) luôn benchmark trên server chính thức kèm CV, ablation, code release; (ii) bao phủ paradigm hiện đại (Transformer, SAM-Med); (iii) bàn rõ \textbf{NLP y tế tiếng Việt low-resource} và \textbf{XAI/uncertainty} cho high-risk theo EU AI Act / FDA.

\textbf{Câu 41. Image processing hỗ trợ chẩn đoán thế nào?}\\
Image processing biến scan thô thành \textbf{biomarker định lượng, có thể tái lập}, bổ sung cho mắt người. (1) \textbf{Phát hiện} — CADe đánh dấu tổn thương (nốt phổi, vi vôi hóa). (2) \textbf{Segmentation} — cho thể tích, hình dáng, tốc độ phát triển (RECIST). (3) \textbf{Classification} — lành/ác, độ ác. (4) \textbf{Registration / dọc thời gian} — cùng bệnh nhân nhiều thời điểm, định lượng thay đổi. (5) \textbf{Radiomics} — hàng trăm đặc trưng texture/hình dáng vào model sống còn. (6) \textbf{Hỗ trợ quyết định} — fuse imaging với EHR, lab, gen cho điều trị cá nhân hoá. Bác sĩ \textbf{vẫn ở trong vòng}; AI là người đọc thứ hai, không thay thế.

\textbf{Câu 42. Pipeline CDSS end-to-end.}\\
(1) \textbf{Thu nhận dữ liệu}: ảnh DICOM từ PACS, EHR (HL7/FHIR), lab, đơn thuốc, giải phẫu bệnh, gen. (2) \textbf{Tiền xử lý}: registration, normalization, NLP trên văn bản tiếng Việt. (3) \textbf{Trích đặc trưng}: biomarker ảnh (segment $\to$ thể tích, radiomics) + đặc trưng cấu trúc. (4) \textbf{Inference}: model đa modality (imaging-CNN $+$ tabular MLP $+$ text BERT) ra điểm rủi ro hoặc chẩn đoán phân biệt. (5) \textbf{Lớp XAI}: Grad-CAM heatmap, SHAP cho tabular, trích dẫn cho text. (6) \textbf{Uncertainty}: tập conformal, xác suất hiệu chuẩn. (7) \textbf{Khuyến nghị}: chẩn đoán xếp hạng + xét nghiệm tiếp theo, theo CPG. (8) \textbf{Giao diện bác sĩ}: tích hợp EMR, một-cú-click chấp nhận / phản hồi, có log audit. (9) \textbf{Vòng phản hồi}: kết quả thực được ghi nhận, dùng retrain.

\textbf{Câu 43. XAI trong AI y tế.}\\
\textbf{XAI} (eXplainable AI) làm cho model black-box dễ hiểu; bắt buộc trong AI y tế để tin cậy, debug, tuân thủ. Phương pháp chính: (1) \textbf{Grad-CAM} — bản đồ kích hoạt theo gradient; heatmap "vùng mạng nhìn" trên CT/MRI. (2) \textbf{Integrated Gradients} / \textbf{SmoothGrad} — quy gán pixel. (3) \textbf{SHAP / LIME} — model-agnostic, phổ biến cho EHR tabular. (4) \textbf{Attention rollout} cho ViT. (5) \textbf{Conformal prediction} — không cần phân phối, cho \textbf{tập dự đoán có coverage hiệu chuẩn} (ví dụ: tập 90\% "chứa chẩn đoán đúng"). (6) \textbf{Dempster--Shafer} — kết hợp bằng chứng đa nguồn. \textbf{EU AI Act} bắt buộc cho CDSS high-risk.

\textbf{Câu 44. Khung pháp lý cho AI y tế.}\\
\textbf{Mỹ}: \textbf{FDA 510(k)} (tương đương predicate device), \textbf{De Novo}, \textbf{PMA} (rủi ro cao nhất), guidance \textbf{Software as a Medical Device (SaMD)}, và mới có \textbf{Predetermined Change Control Plan (PCCP)} cho cập nhật AI. \textbf{Châu Âu}: \textbf{CE marking theo MDR 2017/745}, phân lớp I/IIa/IIb/III. \textbf{EU AI Act} (2024) xếp đa số CDSS là \textbf{high-risk}, bắt buộc quản lý rủi ro, quản trị dữ liệu, minh bạch, giám sát con người, độ chính xác, tính bền, an ninh mạng, giám sát hậu thị trường. \textbf{ISO 13485} (QMS), \textbf{IEC 62304} (vòng đời phần mềm), \textbf{ISO 14971} (rủi ro). Việt Nam dựa Thông tư 30/2015 BYT và đang dần áp dụng các ISO.

\textbf{Câu 45. Thách thức triển khai AI y tế ở Việt Nam.}\\
(1) \textbf{Khan hiếm dữ liệu}: ít corpus y tế tiếng Việt có nhãn công khai; PACS bị khoá theo bệnh viện. (2) \textbf{Ngôn ngữ}: bệnh án trộn \textbf{Việt-Anh-Latin} (ví dụ "U não thể glioblastoma multiforme grade IV"); chưa có KG y tế tiếng Việt phổ biến; ánh xạ ICD-10 không nhất quán. (3) \textbf{Hạ tầng}: PACS cũ, EMR phân mảnh, GPU thiếu ở bệnh viện tỉnh. (4) \textbf{Domain shift}: model train ở nước ngoài (BraTS, ImageNet) kém trên máy quét VN và dân số VN. (5) \textbf{Pháp lý}: chưa có hành lang AI y tế rõ ràng; dựa Thông tư 30/2015 \& ISO. (6) \textbf{Riêng tư}: PDPL 2023 nhưng tiêu chuẩn ẩn danh hóa còn yếu. (7) \textbf{Niềm tin}: bác sĩ cần XAI tiếng Việt. (8) \textbf{Nhân lực}: thiếu chuyên gia tin học y tế. Giải pháp: federated learning xuyên bệnh viện, KG y tế tiếng Việt, fine-tune foundation model bản địa (PhoBERT-Med, ViMedSAM).

\switchcolumn*

\end{paracol}
